\documentclass[11pt,a4paper]{article}
\usepackage[margin=2.2cm]{geometry}
% Preámbulo común de los documentos (article). Compatible con Overleaf (pdfLaTeX, sin -shell-escape).
\usepackage[utf8]{inputenc}
\usepackage[T1]{fontenc}
\usepackage[spanish,es-noshorthands]{babel}
\usepackage{lmodern}
\usepackage{microtype}
\usepackage{amsmath,amssymb}
\usepackage{booktabs}
\usepackage{array}
\usepackage{tabularx}
\usepackage{xcolor}
\usepackage{listings}
\usepackage{enumitem}
\usepackage{caption}
\usepackage[hidelinks]{hyperref}

\captionsetup{font=small,labelfont=bf}
\setlist{itemsep=1pt,topsep=3pt}

\lstset{
  basicstyle=\ttfamily\footnotesize,
  columns=fullflexible,
  keepspaces=true,
  frame=single,
  framerule=0.3pt,
  rulecolor=\color{gray!60},
  backgroundcolor=\color{gray!6},
  breaklines=true,
  inputencoding=utf8,
}

\newcommand{\rel}[1]{\textbf{#1}}
\newcommand{\code}[1]{\texttt{#1}}


\title{Guía de etiquetado de relevancia (escala 0--3)\\[2pt]
\large Listwise Learning to Rank sobre consultas LOINC}
\author{Machine Learned Ranking --- MUIA, UPM}
\date{Noviembre 2026}

\begin{document}
\maketitle

\section{Objetivo}

Asignar a cada par consulta--documento una etiqueta de relevancia ordinal en $\{0,1,2,3\}$, de forma
\textbf{reproducible y justificada}. Los documentos son los 67 códigos LOINC del libro
\code{loinc\_dataset-v2.xlsx}, que aparecen (en distinto orden) en las tres hojas, una por consulta:
\emph{glucose in blood}, \emph{bilirubin in plasma} y \emph{white blood cells count}. En total se juzgan
$3 \times 67 = 201$ pares.

El orden de las filas en las hojas \textbf{no} se interpreta como etiqueta. La misma prueba puede ser
relevante para una consulta e irrelevante para otra, por lo que cada par se juzga de forma independiente.

La regla está implementada en \code{src/queries.py} y se aplica con \code{src/label.py}, que genera
\code{data/relevance\_labels.csv} con una justificación (\code{rationale}) por cada par.

\section{Escala}

\begin{center}
\begin{tabularx}{\linewidth}{c l X}
\toprule
Etiqueta & Nombre & Significado \\
\midrule
\rel{3} & Muy relevante & Mide directamente el analito pedido, en la muestra pedida y con una medida adecuada. \\
\rel{2} & Relevante con un matiz & Satisface la consulta salvo por \emph{una} desviación: muestra compatible pero distinta, fracción del analito o condición especial. \\
\rel{1} & Relación parcial & Comparte el analito, pero no satisface la consulta (muestra incompatible, medida cualitativa, subpoblación expresada como porcentaje\ldots). \\
\rel{0} & Irrelevante & El analito es otro, aunque comparta palabras con la consulta (p.\,ej.\ \emph{Blood product type}). \\
\bottomrule
\end{tabularx}
\end{center}

\section{Las tres dimensiones del juicio}

Cada par se evalúa en tres dimensiones, que se leen de las columnas LOINC del documento:

\begin{description}
\item[Analito] (columna \code{component}).
  \emph{Exacto}: es el concepto pedido.
  \emph{Relacionado}: es una fracción o subpoblación del concepto pedido.
  \emph{Ninguno}: es otro analito.
\item[Muestra] (columna \code{system}).
  \emph{Exacta}: es la muestra pedida, o un sistema LOINC que la incluye.
  \emph{Compatible}: es una muestra distinta pero clínicamente intercambiable.
  \emph{Incompatible}: cualquier otra muestra (orina, líquidos biológicos, muestra no especificada \code{XXX}\ldots).
\item[Medida/condición] (columna \code{property} y sufijo \code{\^{}} de \code{component}).
  \emph{Adecuada}: la propiedad es la que pide la consulta y no hay condiciones especiales.
  \emph{Desviada}: la propiedad no es la pedida (cualitativa \code{PrThr}/\code{ACnc}, fracción \code{NFr} cuando se pide un recuento\ldots),
  o el componente lleva una condición tras \code{\^{}} (ayuno, post-carga de glucosa, corrección por albúmina\ldots).
\end{description}

\section{Regla de cálculo}

\begin{enumerate}
\item Si el analito es \emph{ninguno}, la etiqueta es \rel{0}.
\item Base: \rel{3} si el analito es \emph{exacto}; \rel{2} si es \emph{relacionado}.
\item Se resta 1 si la muestra es \emph{compatible} y 2 si es \emph{incompatible}.
\item Se resta 1 si la medida/condición está \emph{desviada}.
\item Si el analito es exacto o relacionado, la etiqueta mínima es \rel{1}: compartir analito nunca equivale a irrelevante.
\end{enumerate}

\section{Criterios por consulta}

\begin{center}
\small
\begin{tabularx}{\linewidth}{>{\raggedright}p{2.6cm} >{\raggedright}X >{\raggedright}p{2.6cm} >{\raggedright}p{2.2cm} >{\raggedright\arraybackslash}p{1.9cm}}
\toprule
Consulta & Analito exacto / relacionado & Muestra exacta & Muestra compatible & Propiedad adecuada \\
\midrule
glucose in blood & \code{Glucose} (con o sin condición \code{\^{}}) / --- & \code{Bld}, \code{BldC}, \code{Ser/Plas/Bld} & \code{Ser/Plas}, \code{Ser}, \code{Plas} & \code{MCnc}, \code{SCnc}, \code{MSCnc} \\
bilirubin in plasma & \code{Bilirubin} (total) / fracciones directa (glucuronidada) e indirecta (no glucuronidada) & \code{Plas}, \code{Ser/Plas} & \code{Ser}, \code{Bld} & \code{MCnc}, \code{SCnc}, \code{MSCnc} \\
white blood cells count & \code{Leukocytes} / subpoblaciones: \code{Lymphocytes}, \code{Monocytes}, neutrófilos\ldots & \code{Bld} & --- & \code{NCnc} (recuento por volumen) \\
\bottomrule
\end{tabularx}
\end{center}

\subsection*{Decisiones de diseño}

\begin{itemize}
\item \textbf{Suero/plasma frente a sangre.} \code{Ser/Plas} se considera una muestra \emph{compatible pero distinta} de
  \code{Bld}: el suero y el plasma son fracciones de la sangre y la glucemia se mide habitualmente en ellos, pero no son
  sangre total. Por eso \emph{Glucose [Moles/volume] in Serum or Plasma} recibe \rel{2} para \emph{glucose in blood}. De forma simétrica,
  \code{Bld} es compatible pero distinta para \emph{bilirubin in plasma}. \code{Ser/Plas/Bld} incluye explícitamente la sangre, así que es exacta.
\item \textbf{Bilirrubina total frente a fracciones.} Una consulta sin fracción (\emph{bilirubin}) se interpreta como bilirrubina total.
  La directa y la indirecta son fracciones (analito relacionado, base 2).
\item \textbf{WBC = white blood cells = leukocytes.} El recuento total es \code{Leukocytes [\#/volume]}. Los recuentos de
  subpoblaciones (linfocitos, monocitos) son analito relacionado (\rel{2}). Un \emph{porcentaje} de una subpoblación
  (\code{Lymphocytes/100 leukocytes}, \code{NFr}) no es un recuento, por lo que baja a \rel{1}.
\item \textbf{Condiciones especiales.} La glucosa en ayunas o tras una sobrecarga oral mide el mismo analito, pero responde a una
  pregunta clínica distinta (prueba funcional). Por eso cuenta como medida desviada.
\item \textbf{Coincidencias léxicas engañosas.} \emph{Blood product type}, \emph{Blood group antibody screen}, \emph{ABO group in Blood}\ldots{}
  contienen \emph{blood}, pero su analito no es glucosa ni leucocitos, así que reciben \rel{0}.
  \emph{Calcium bilirubinate/Total in Stone} describe la composición de un cálculo biliar, no la bilirrubina circulante, así que recibe \rel{0}.
\item \textbf{Empates.} Documentos con la misma etiqueta pueden empatar: no se impone un orden artificial entre ellos.
\end{itemize}

\section{Resultado}

\begin{table}[h]
\centering
% Generado por src/export_tables.py — no editar a mano
\begin{tabular}{lrrrrr}
\toprule
Consulta & 3 & 2 & 1 & 0 & Total \\
\midrule
glucose in blood & 1 & 1 & 4 & 61 & 67 \\
bilirubin in plasma & 1 & 3 & 2 & 61 & 67 \\
white blood cells count & 1 & 2 & 1 & 63 & 67 \\
\midrule
calcium in serum$^\dagger$ & 1 & 2 & 1 & 63 & 67 \\
cholesterol in plasma$^\dagger$ & 1 & 1 & 0 & 65 & 67 \\
chloride in blood$^\dagger$ & 1 & 1 & 0 & 65 & 67 \\
\bottomrule
\end{tabular}

\caption{Distribución de etiquetas por consulta. $^\dagger$Consultas adicionales (parte opcional),
etiquetadas con la misma regla sobre los mismos 67 documentos.}
\end{table}

\begin{table}[h]
\centering
\footnotesize
% Generado por src/export_tables.py — no editar a mano
\begin{tabular}{llp{6.2cm}llc}
\toprule
Consulta & LOINC & long\_common\_name & system & property & Rel. \\
\midrule
glucose in blood & 74774-1 & Glucose [Mass/volume] in Serum, Plasma or Blood & Ser/Plas/Bld & MCnc & 3 \\
 & 14749-6 & Glucose [Moles/volume] in Serum or Plasma & Ser/Plas & SCnc & 2 \\
 & 15076-3 & Glucose [Moles/volume] in Urine & Urine & SCnc & 1 \\
 & 14747-0 & Glucose [Moles/volume] in Pleural fluid & Plr fld & SCnc & 1 \\
 & 14764-5 & Glucose [Moles/volume] in Serum or Plasma --3 hours post 100 g glucose PO & Ser/Plas & SCnc & 1 \\
 & 35184-1 & Fasting glucose [Mass or Moles/volume] in Serum or Plasma & Ser/Plas & MSCnc & 1 \\
\midrule
bilirubin in plasma & 1975-2 & Bilirubin.total [Mass/volume] in Serum or Plasma & Ser/Plas & MCnc & 3 \\
 & 1968-7 & Bilirubin.direct [Mass/volume] in Serum or Plasma & Ser/Plas & MCnc & 2 \\
 & 1971-1 & Bilirubin.indirect [Mass/volume] in Serum or Plasma & Ser/Plas & MCnc & 2 \\
 & 35192-4 & Bilirubin.indirect [Mass or Moles/volume] in Serum or Plasma & Ser/Plas & MSCnc & 2 \\
 & 14423-8 & Bilirubin.total [Mass/volume] in Synovial fluid & Synv fld & MCnc & 1 \\
 & 33870-7 & Bilirubin.total [Presence] in Unspecified specimen & XXX & PrThr & 1 \\
\midrule
white blood cells count & 26464-8 & Leukocytes [\#/volume] in Blood & Bld & NCnc & 3 \\
 & 26474-7 & Lymphocytes [\#/volume] in Blood & Bld & NCnc & 2 \\
 & 26484-6 & Monocytes [\#/volume] in Blood & Bld & NCnc & 2 \\
 & 26478-8 & Lymphocytes/100 leukocytes in Blood & Bld & NFr & 1 \\
\bottomrule
\end{tabular}

\caption{Todos los documentos con relevancia mayor que 0 en las consultas oficiales. El resto de documentos de cada
consulta (61--63) son \rel{0} porque su analito es distinto.}
\end{table}

\section{Revisión manual y overrides}

La regla es determinista, pero las 201 justificaciones se revisan a mano. Si un juicio de la regla se considera
incorrecto, se añade una fila a \code{data/overrides.csv} (\code{qid, loinc\_num, relevance, justification}).
\code{src/label.py} aplica el override y deja constancia en el \code{rationale}
(\code{OVERRIDE a -> b: motivo}). En la versión actual \textbf{no hay overrides}. Cualquier override futuro
debe justificarse en esta sección.

\paragraph{Diferencias con la versión anterior (escala 0--2).} Respecto al \code{data/dataset\_v0.csv} inicial,
además de pasar a cuatro niveles, cambia el juicio en estos casos:
\begin{itemize}
\item \emph{Blood product type}: era 1 para glucosa y pasa a \rel{0} (otro analito).
\item \emph{Fasting glucose}: era 0 y pasa a \rel{1} (mismo analito, con condición y muestra compatible).
\item \emph{Glucose in Serum or Plasma}: empataba con la glucosa en orina (1) y ahora queda por encima (\rel{2} frente a \rel{1}).
\item \emph{Lymphocytes/100 leukocytes}: empataba con el recuento total de leucocitos (2) y ahora es \rel{1} (porcentaje, no recuento).
\item \emph{Monocytes [\#/volume]}: estaba por debajo de \emph{Lymphocytes [\#/volume]} y ahora ambos son \rel{2} (subpoblaciones).
\item Se corrige la errata \emph{bilirrubin} $\rightarrow$ \emph{bilirubin} en el texto de la consulta.
\end{itemize}

\end{document}
